% Options for packages loaded elsewhere
\PassOptionsToPackage{unicode}{hyperref}
\PassOptionsToPackage{hyphens}{url}
%
\documentclass[
]{book}
\usepackage{amsmath,amssymb}
\usepackage{lmodern}
\usepackage{ifxetex,ifluatex}
\ifnum 0\ifxetex 1\fi\ifluatex 1\fi=0 % if pdftex
  \usepackage[T1]{fontenc}
  \usepackage[utf8]{inputenc}
  \usepackage{textcomp} % provide euro and other symbols
\else % if luatex or xetex
  \usepackage{unicode-math}
  \defaultfontfeatures{Scale=MatchLowercase}
  \defaultfontfeatures[\rmfamily]{Ligatures=TeX,Scale=1}
\fi
% Use upquote if available, for straight quotes in verbatim environments
\IfFileExists{upquote.sty}{\usepackage{upquote}}{}
\IfFileExists{microtype.sty}{% use microtype if available
  \usepackage[]{microtype}
  \UseMicrotypeSet[protrusion]{basicmath} % disable protrusion for tt fonts
}{}
\makeatletter
\@ifundefined{KOMAClassName}{% if non-KOMA class
  \IfFileExists{parskip.sty}{%
    \usepackage{parskip}
  }{% else
    \setlength{\parindent}{0pt}
    \setlength{\parskip}{6pt plus 2pt minus 1pt}}
}{% if KOMA class
  \KOMAoptions{parskip=half}}
\makeatother
\usepackage{xcolor}
\IfFileExists{xurl.sty}{\usepackage{xurl}}{} % add URL line breaks if available
\IfFileExists{bookmark.sty}{\usepackage{bookmark}}{\usepackage{hyperref}}
\hypersetup{
  pdftitle={Probability, R and Computational Finance},
  pdfauthor={Martin Summer and Branko Urosevic},
  hidelinks,
  pdfcreator={LaTeX via pandoc}}
\urlstyle{same} % disable monospaced font for URLs
\usepackage{longtable,booktabs,array}
\usepackage{calc} % for calculating minipage widths
% Correct order of tables after \paragraph or \subparagraph
\usepackage{etoolbox}
\makeatletter
\patchcmd\longtable{\par}{\if@noskipsec\mbox{}\fi\par}{}{}
\makeatother
% Allow footnotes in longtable head/foot
\IfFileExists{footnotehyper.sty}{\usepackage{footnotehyper}}{\usepackage{footnote}}
\makesavenoteenv{longtable}
\usepackage{graphicx}
\makeatletter
\def\maxwidth{\ifdim\Gin@nat@width>\linewidth\linewidth\else\Gin@nat@width\fi}
\def\maxheight{\ifdim\Gin@nat@height>\textheight\textheight\else\Gin@nat@height\fi}
\makeatother
% Scale images if necessary, so that they will not overflow the page
% margins by default, and it is still possible to overwrite the defaults
% using explicit options in \includegraphics[width, height, ...]{}
\setkeys{Gin}{width=\maxwidth,height=\maxheight,keepaspectratio}
% Set default figure placement to htbp
\makeatletter
\def\fps@figure{htbp}
\makeatother
\setlength{\emergencystretch}{3em} % prevent overfull lines
\providecommand{\tightlist}{%
  \setlength{\itemsep}{0pt}\setlength{\parskip}{0pt}}
\setcounter{secnumdepth}{5}
\usepackage{booktabs}
\ifluatex
  \usepackage{selnolig}  % disable illegal ligatures
\fi
\usepackage[]{natbib}
\bibliographystyle{apalike}

\title{Probability, R and Computational Finance}
\author{Martin Summer and Branko Urosevic}
\date{2021-07-13}

\begin{document}
\maketitle

{
\setcounter{tocdepth}{1}
\tableofcontents
}
\hypertarget{introduction}{%
\chapter{Introduction}\label{introduction}}

\begin{quote}
``What I hear, I forget; What I see, I remember; What I do, I understand''

\hfill --- Confucius, 551-479 BC
\end{quote}

Our lectures are going to teach you basic concepts of probability and
their applications in computational finance. Our approach is based on
intuition and on understanding concepts by building them with your own hands.

But how can we build abstract concepts such as probability, random phenomena and chance by
our own hands? How can we ``touch'' such things as ideas?
While probability is a mathematical theory, it gains practical value and an intuitive meaning in
connection with real or conceptual experiments such as, tossing a coin, rolling two dice, playing
roulette, observing the change in the index value of stocks, bonds or in exchange rates, selecting
a random sample of financial transactions in a payment system and observing
the number of transactions with a value below a certain amount.

Many of these experiments we can nowadays do not only physically -- by actually
rolling a pair of dice, for example -- but we can use the computer to build a
virtual dice and roll it on the computer. We
can build and simulate a huge variety of random phenomena. We can - for instance -
implement models of random fluctuations of asset prices. We can model financial risks
and contemplate possible future scenarios through simulation.

We can approach our understanding of abstract concepts
by building them with our own hands on the computer. This is the approach to teaching you
probability in this course.

There are many ways to build virtual objects and to run simulations to manipulate
them. For this we will need a programming language. The language we choose for this
course is R and the integrated development environment RStudio. This is one of the main
languages used in data analysis, statistics and data science and is widely used in industry and
academia. It will be our tool to do probability in this course.

\hypertarget{downloading-and-installing-r}{%
\section{Downloading and installing R}\label{downloading-and-installing-r}}

So lets start by downloading and installing R first. R is an open source project maintained
by an international team of developers. The software is made available through a website called
the comprehensive R archive network (\url{http://cran.r-project.org}).

At the top of this website, in a box named ``Download and install R'' you will find three
links for downloading R. Choose the link that describes your operating system, Windows, Mac or
Linux. These links will lead you to the necessary information you need to install a current
version of R. The easiest install option is to install R from precompiled binaries. There is also
the option to built R from source on all operationg systems if you have the tools and the
expertise to do so. R also comes both in a 32-bit and a 64-bit version. It does not make
a substantial differecne which version you use. 64-bit versions
can handle larger files and datasets with fewer memory management problems.

\hypertarget{downloading-and-installing-rstudio}{%
\section{Downloading and installing RStudio}\label{downloading-and-installing-rstudio}}

RStudio is an application that helps you write and develop R code. It makes using
R much easier fro you than using it in isolation. The interface of RStudio looks the
same across all operating systems.

You can download RStudio for free from \url{https://www.rstudio.com/products/rstudio/}.
Select the box RStudio Desktop and follow the download instructions. RStudio Desktop is free.
Note that you need to have a version of R installed to use RStudio.

If you have successfully installed R and RStudio, we are ready to start.

\hypertarget{empirical-background-and-some-terminology}{%
\chapter{Empirical background and some terminology}\label{empirical-background-and-some-terminology}}

If you chose to study Finance it should be natural for you to study probability.
After all, Finance studies the allocation and pricing of money across time. You can
safe money to use it later for purchasing goods and services you can borrow money to
invest today and repay in the future from the revenues of your project. But the
future is uncertain and we do not know today how it will look like. In Finance we unavoidably
have to deal with risk and uncertainty. The most fundamental concept we have available
to deal with risk and uncertainty and to think about and analyze it is probability. Therefore
every serious student of Finance has to study probability early on.

Probability is a mathematical theory. This theory gets practical value and intuitive meaning
in connection with real or conceptual experiments such as rolling a dice, tossing a coin or
selecting a random sample of stocks and observing their most extreme price drop over a
given period of time or observing the frequency of car accidents etc. All of these
descriptions are, however, rather vague. We have to be more precise what we mean by an
experiment or observation.

When we use the term (random) \emph{experiment} we mean a process leading to an uncertain
outcome. To make this notion precise and useful we need to agree what we mean by
uncertain outcomes.

Take the simple example of tossing a coin. In a practical situations the outcome
of a coin toss can be that the coin lands on head or tail. But it could in principle
also fall on an edge, stand and roll away. Still when we think about the
experiment of tossing a coin in the theory of probability we agree \emph{on the outset} that
heads and tail are the only possible outcomes of this experiment. By such an
agreement we have pinned down the possible outcomes.

The collection of all possible outcomes of an experiment is calles the \emph{sample space}.
While the sample space of the experiment of tossing a coin is an idealization it is
exactly this idealization which simplifies the theory without affecting its applicability.

Idealizations of this kind are standard in probability and are widely used also in Finance. It
is - for example - impossible to measure how high the price of a stock might become. Is there
a maximal price beyond which the price can not possibly be higher? Many of us would
hesitate to claim that the price might rise without bound. Yet many models in applied
Finance are based on such an assumption. The models allow arbitrary price hikes but with arbitrary
small probability as the price gets higher and higher. Practically it does not make sense
to believe that a security price can become arbitratily high. The use of arbitrarily small
probabilities in a financial model might seem absurd but it does no practical harm and makes
the model simple and convenient to use. Moreover, if we seriously introduced an upper bound
on a security price at \(x\) it would be also awkward to assume that it is impossible that it
could be just a cent higher, an assumption equally unappealing as assuming it can get in
principle arbitrarily high.

When we talk of a \emph{basic outcome} we mean the possible outcome of a random experiment. An \emph{event} is
a subset of basic outcomes. Any event which consists of a single outcome in the sample space
is called a \emph{simple event}.

Probability is a measure of how likely an event of an experiment is. Thus when we
talk about probability in a precise and meaningful sense we can only so in relation to
a given sample space or to a certain conceptual experiment. Thus before we can begin to seriously
talk about probability we need to develop a good understanding of experiments, sample
spaces, basic outcomes and events. In the spirit of this course we will develop this
understanding by building these concepts using R. This will at the same time
get us started with using the R language as a tool which will allow us to ``touch''
probability concepts. As we proceed through our course we will develop a broader and deeper
understanding of this language together with the probability concepts.

\hypertarget{building-experiments-sample-spaces-and-events-using-r}{%
\chapter{Building experiments, sample spaces and events using R}\label{building-experiments-sample-spaces-and-events-using-r}}

\hypertarget{literature}{%
\chapter{Literature}\label{literature}}

Here is a review of existing methods.

\hypertarget{methods}{%
\chapter{Methods}\label{methods}}

We describe our methods in this chapter.

\hypertarget{applications}{%
\chapter{Applications}\label{applications}}

Some \emph{significant} applications are demonstrated in this chapter.

\hypertarget{example-one}{%
\section{Example one}\label{example-one}}

\hypertarget{example-two}{%
\section{Example two}\label{example-two}}

\hypertarget{final-words}{%
\chapter{Final Words}\label{final-words}}

We have finished a nice book.

\hypertarget{appendix}{%
\chapter{Appendix}\label{appendix}}

\hypertarget{installing-r}{%
\section{Installing R}\label{installing-r}}

\hypertarget{installing-rstudio}{%
\section{Installing RStudio}\label{installing-rstudio}}

  \bibliography{book.bib,packages.bib}

\end{document}
